% dataset: math-formula-atlas

% page_slug: exponent-fractional-html

% page_title: Fractional Exponents

% page_url: https://www.mathsisfun.com/algebra/exponent-fractional.html

% subject_slug: algebra

% subject_path: /algebra/

% formula_count: 11

% formula_record_start
% formula_id: exponent-fractional-html-001
% index_global: 43
% index_in_page: 1
% section: Try Another Fraction
% subsection: none
% formula_name: 
% source: formula-text-fallback
\[Let us try that again, but with an exponent of one-quarter (\\frac{1}{4}):\]
% formula_record_end

% formula_record_start
% formula_id: exponent-fractional-html-002
% index_global: 44
% index_in_page: 2
% section: General Rule
% subsection: none
% formula_name: 
% source: formula-text-fallback
\[x^{1/n} = The n-th Root of x\]
% formula_record_end

% formula_record_start
% formula_id: exponent-fractional-html-003
% index_global: 45
% index_in_page: 3
% section: What About More Complicated Fractions?
% subsection: none
% formula_name: 
% source: formula-text-fallback
\[What about a fractional exponent like 4^{\\frac{3}{2}} ?\]
% formula_record_end

% formula_record_start
% formula_id: exponent-fractional-html-004
% index_global: 46
% index_in_page: 4
% section: What About More Complicated Fractions?
% subsection: none
% formula_name: 
% source: formula-text-fallback
\[That is really saying to do a cube (3) and a square root (\\frac{1}{2}), in any order.\]
% formula_record_end

% formula_record_start
% formula_id: exponent-fractional-html-005
% index_global: 47
% index_in_page: 5
% section: What About More Complicated Fractions?
% subsection: none
% formula_name: 
% source: formula-text-fallback
\[a fraction (1/n) part\]
% formula_record_end

% formula_record_start
% formula_id: exponent-fractional-html-006
% index_global: 48
% index_in_page: 6
% section: What About More Complicated Fractions?
% subsection: none
% formula_name: 
% source: formula-text-fallback
\[So, because m/n = m \times (1/n) we can do this:\]
% formula_record_end

% formula_record_start
% formula_id: exponent-fractional-html-007
% index_global: 49
% index_in_page: 7
% section: What About More Complicated Fractions?
% subsection: none
% formula_name: 
% source: formula-text-fallback
\[x^{m/n} = x^{(m \times 1/n)} = (x^{m})^{1/n} = n\sqrt{}x^{m}\]
% formula_record_end

% formula_record_start
% formula_id: exponent-fractional-html-008
% index_global: 50
% index_in_page: 8
% section: What About More Complicated Fractions?
% subsection: none
% formula_name: 
% source: formula-text-fallback
\[The order does not matter, so it also works for m/n = (1/n) \times m:\]
% formula_record_end

% formula_record_start
% formula_id: exponent-fractional-html-009
% index_global: 51
% index_in_page: 9
% section: What About More Complicated Fractions?
% subsection: none
% formula_name: 
% source: formula-text-fallback
\[x^{m/n} = x^{(1/n \times m)} = (x^{1/n})^{m} = (n\sqrt{}x )^{m}\]
% formula_record_end

% formula_record_start
% formula_id: exponent-fractional-html-010
% index_global: 52
% index_in_page: 10
% section: Now ... Play With The Graph!
% subsection: none
% formula_name: 
% source: formula-text-fallback
\[Start with m=1 and n=1, then slowly increase n so that you can see \\frac{1}{2}, \\frac{1}{3} and \\frac{1}{4}\]
% formula_record_end

% formula_record_start
% formula_id: exponent-fractional-html-011
% index_global: 53
% index_in_page: 11
% section: Now ... Play With The Graph!
% subsection: none
% formula_name: 
% source: formula-text-fallback
\[Then try m=2 and slide n up and down to see fractions like \\frac{2}{3} and so on\]
% formula_record_end